% paper/sections/09b_implementation.tex
\section{Rhie--Chow・Balanced-Force：実装整合性と検証手順}
\label{sec:impl_collocate}

第~\ref{sec:algorithm}章（完全アルゴリズムフロー）で示した時間積分ループでは，
Rhie--Chow 補正と Balanced-Force 離散化が正しく協調していることが安定動作の前提となる．
本章では，第~\ref{sec:collocate}章で導いた Rhie--Chow 補正（§\ref{sec:rhie_chow}）および
Balanced-Force 離散化条件（§\ref{sec:balanced_force}）を本コードで実現するにあたっての
実装上の整合性要件と，静止液滴テストによる検証手順を述べる．

% ─────────────────────────────────────────────────────────────────────────────
\subsection{コード実装との整合性：演算子の統一確認}
\label{sec:impl_bf_consistency}
% ─────────────────────────────────────────────────────────────────────────────

§\ref{sec:balanced_force} の議論より，
式~\eqref{eq:bf_operator_mismatch} の離散演算子不整合をゼロにするには，
圧力勾配・表面張力・速度補正の 3 箇所すべてに同一の $D^{(1)}_\mathrm{CCD}$ を用いなければならない．
本稿のコードでは以下の対応関係で実装している：

\begin{enumerate}
  \item \texttt{SurfaceTensionTerm.compute()}:
        $\bm{f}_\sigma = (\kappa/We)\,D^{(1)}_\mathrm{CCD}\psi$
  \item \texttt{Predictor.compute()} IPC 項:
        $-D^{(1)}_\mathrm{CCD} p^n$（前時刻圧力勾配）
  \item \texttt{VelocityCorrector.correct()}:
        $-({\Delta t}/{\rho})\,D^{(1)}_\mathrm{CCD}(\delta p)$（速度補正）
\end{enumerate}

これにより，式~\eqref{eq:bf_operator_mismatch} の不整合はゼロに設計されている．

% ─────────────────────────────────────────────────────────────────────────────
\subsection{Balanced-Force RC：実装状況}
\label{sec:impl_bf_rc}
% ─────────────────────────────────────────────────────────────────────────────

\texttt{rhie\_chow.py} の \texttt{face\_velocity\_divergence} は，
式~\eqref{eq:rc-face}（圧力補正のみ，デフォルト）と
式~\eqref{eq:rc-face-balanced}（BF-RC 拡張，オプション）の両形式を実装している．
\texttt{kappa}，\texttt{psi}，\texttt{we} を引数として渡すと
式~\eqref{eq:rc-face-balanced} が有効化される．
定量的検証は \texttt{experiments/balanced\_force\_rc\_benchmark.py}
（§\ref{sec:bench_droplet} 参照）で実施できる．

% ─────────────────────────────────────────────────────────────────────────────
\subsection{Balanced-Force 実装の確認方法：静止液滴テスト}
\label{sec:impl_bf_verify}
% ─────────────────────────────────────────────────────────────────────────────

\begin{tcolorbox}[algbox, title={静止液滴テスト：Balanced-Force の検証手順}]
\begin{enumerate}
  \item \textbf{初期条件：} 計算領域中央に半径 $R$ の円形液滴を置く．
        速度 $\bm{u} = \bm{0}$，圧力は Young-Laplace 解
        $p^0 = -\sigma\kappa_0 = \sigma/R$（液相内，$\kappa_0 = -1/R$），
        $p^0 = 0$（気相）．

  \item \textbf{時間発展：} 数十タイムステップ計算する．
        Balanced-Force が正しく実装されていれば，速度は CSF モデル誤差が律速する
        $O(\varepsilon^2) \approx O(h^2)$ レベルに留まる
        （例：$N=64$ で $\|\bm{u}\|_{L^\infty} \lesssim 10^{-3}$〜$10^{-4}$ 程度）．
        CSF が界面を有限厚さ $\varepsilon \sim h$ で近似する限り，
        \textbf{モデル誤差 $O(h^2)$} が寄生流れの下限を決めるため，
        機械精度への収束は物理的に不可能である（付録~\ref{app:balanced_force_taylor}）．

  \item \textbf{収束確認：} 格子幅 $h$ を変えながら最大速度振幅
        \[
          U_{\mathrm{max}}(h) = \|\bm{u}\|_{L^\infty}
        \]
        を測定し，$U_{\mathrm{max}}(h) = O(h^q)$（$q \approx 2$）の収束を確認する．
        CSF モデルの界面厚さ $\varepsilon \sim O(h)$ に起因するモデル誤差 $O(\varepsilon^2) \approx O(h^2)$
        が律速するため，Balanced-Force + CCD を正しく実装しても収束次数は $q \approx 2$ となる
        （付録~\ref{app:balanced_force_taylor}，§3）．
        Balanced-Force の効果は収束\textbf{次数}の改善ではなく，収束\textbf{係数}の大幅な低減
        ——例えば $U_{\mathrm{max}}$ が $O(10^{-1})$〜$O(1)$（実装ミス）から
        $O(10^{-3})$〜$O(10^{-4})$（正しい実装）へ——として現れる（§\ref{sec:bench_droplet} 参照）．

  \item \textbf{失敗の兆候：} $U_{\mathrm{max}}$ が $O(1)$〜$O(10^{-2})$ 程度で収束しない場合，
        以下を確認する：
        \begin{itemize}
          \item 圧力勾配と表面張力に同一の差分演算子を使っているか？
          \item Rhie--Chow の密度係数と PPE の係数に同じ調和平均を使っているか？
          \item 曲率計算と表面張力の計算タイミングが整合しているか？
          \item 表面張力支配の問題（$We \ll 1$）では，式~\eqref{eq:rc-face-balanced} の
                表面張力補正項を Rhie--Chow に追加しているか（§\ref{sec:rc_balanced_force} 参照）？
        \end{itemize}
\end{enumerate}
\end{tcolorbox}
